\section{SOLUÇÃO DESENVOLVIDA: WeFIND}

A concepção do WeFIND foi motivada por uma experiência vivenciada pelo autor durante a busca por um animal doméstico desaparecido no círculo familiar. Na ocasião, o acompanhamento da busca dependia de publicações distribuídas em redes sociais e grupos de mensagens, o que dificultava localizar as informações, identificar animais anunciados nas proximidades e acompanhar as atualizações. Essa experiência evidenciou, em uma situação concreta, a necessidade de uma ferramenta capaz de reunir as informações e organizar a colaboração entre as pessoas envolvidas.

A partir dessa motivação, a plataforma WeFIND foi desenvolvida para centralizar informações sobre animais perdidos e encontrados, organizando a participação dos usuários em um ambiente específico para esse fim. Em resposta às dificuldades observadas, a solução reúne recursos de publicação, busca geográfica, comunicação e colaboração comunitária. A experiência pessoal, portanto, constitui a motivação inicial do projeto, enquanto as funcionalidades apresentadas a seguir resultam da análise do problema e dos requisitos definidos ao longo do desenvolvimento.

Para estruturar a solução desenvolvida e contemplar os requisitos levantados, as funcionalidades centrais foram agrupadas em três eixos:

O primeiro eixo compreende a busca espacial georreferenciada, organizando o mapa interativo e o mural de publicações por proximidade territorial e permitindo a delimitação de raios de busca.

O segundo eixo estabelece o registro colaborativo de avistamentos. Por meio desse recurso, os usuários relatam animais observados em vias públicas com fotografias e coordenadas, construindo uma linha do tempo vinculada à publicação correspondente.

O terceiro eixo reúne comunicação, privacidade e divulgação. O sistema oferece um canal de mensagens privado entre os participantes e permite compartilhar cartazes de busca, apoiando a circulação das informações sem expor publicamente o telefone do usuário.


\subsection{Escolha da Plataforma Móvel}

A decisão de desenvolver o WeFIND como um aplicativo para celular fundamentou-se na necessidade de apoiar ações realizadas em deslocamento, nas quais o usuário precisa consultar publicações, registrar avistamentos ou receber alertas diretamente no smartphone.

O uso de recursos geográficos em sistemas de apoio a buscas por animais também é abordado no contexto brasileiro por Souza et al. (2022), que apresentam um algoritmo de recomendação de rotas para orientar a procura por animais desaparecidos. O WeFIND tem um escopo diferente: organiza publicações e relatos de avistamento por localização, sem recomendar o percurso do animal nem afirmar onde ele se encontra. A opção por uma aplicação móvel permite que o usuário consulte publicações e registre informações durante a própria busca, na rua.

A plataforma móvel reúne recursos que apoiam esse contexto. As notificações informam os usuários sobre atualizações relevantes, o acesso à câmera facilita o registro de fotografias e os recursos de localização auxiliam na marcação do ponto do avistamento. A geração de cartazes digitais com código QR viabiliza a divulgação externa impressa e em redes sociais, direcionando os interessados de volta à publicação no aplicativo.


\subsection{Sugestão de Correspondências e Registro de Avistamentos}

O WeFIND adota duas formas distintas e complementares para relacionar informações sobre animais: a correspondência automatizada entre publicações e o registro comunitário de avistamentos. Embora ambas utilizem características físicas e dados de localização, operam com propósitos e fluxos distintos.


\subsubsection{Correspondência entre Publicações}

A correspondência automática analisa duas publicações ativas cadastradas de forma independente: uma do tipo ``Perdido'' e outra do tipo ``Encontrado''. O sistema avalia a compatibilidade de atributos como espécie, porte, sexo, raça, pelagem e idade, cruzando esses dados com a distância geográfica entre os pontos registrados ou, na ausência de coordenadas, com a cidade e o bairro informados. A listagem pode apresentar sugestões a partir de 55 pontos, em um índice heurístico de 0 a 100; as notificações automáticas utilizam o limiar de 60 pontos.

A formulação matemática desse cálculo, a atribuição de pesos a cada atributo morfológico e a aplicação da fórmula de Haversine para a distância geográfica são detalhadas na seção do Capítulo~6 dedicada ao motor de correspondência e no Quadro~4. As interfaces do usuário e as notificações resultantes desse cruzamento são ilustradas nas telas do fluxo de localização e correspondência apresentadas no Capítulo~7.

Essa funcionalidade tem como objetivo aproximar tutores e resgatistas que possam estar tratando do mesmo animal. Ela não altera o status das publicações, não confirma reencontros e não substitui a validação humana: a ferramenta atua como um mecanismo de apoio para que as partes comparem fotos e características, cabendo exclusivamente aos envolvidos a confirmação do caso.


\subsubsection{Registro de Avistamentos}

O registro de avistamento opera a partir do relato de uma pessoa que visualizou um animal na rua, possivelmente associado a uma busca em andamento. O sistema utiliza as coordenadas do informante para consultar publicações ativas em um raio inicial de 5 km, definido no fluxo de busca para restringir os resultados à região relevante.

Após a conferência dos dados e a confirmação do usuário de que o animal visto corresponde a um anúncio específico, o aplicativo registra um novo relato de avistamento vinculado à publicação original. Esse relato preserva descrição, fotografia, coordenadas e identificação do informante em uma linha do tempo cronológica. A publicação principal é atualizada com o último ponto geográfico conhecido, mantendo intacto o histórico de deslocamento do animal. Caso o usuário verifique que se trata de outro animal, nenhuma publicação é alterada, podendo o relato ser descartado ou convertido em um novo cadastro.

Dessa forma, a correspondência automática atua comparando dois anúncios distintos, enquanto o avistamento acrescenta novos relatos a uma busca já acompanhada. Em ambos os fluxos, o sistema apenas organiza as informações, mantendo a autonomia de decisão com os usuários.


\subsection{Categorias e Estados das Publicações}

Para assegurar clareza operacional, o WeFIND estabelece distinção formal entre as categorias de anúncio e os estados de encerramento da publicação.

As categorias iniciais de cadastro compreendem ``Perdido'' (publicação gerada pelo tutor à procura de seu animal) e ``Encontrado'' (anúncio criado por quem localizou um animal ou o observou em via pública).

O termo ``Reencontrado'' representa o encerramento da busca ativa quando o animal retorna ao tutor ou a tutela é confirmada entre as partes.

\subsection{Comprovação de Tutela}

Quando alguém encontra um animal e o anuncia no WeFIND, outra pessoa pode se apresentar como seu tutor. Para reduzir o risco de devolução a quem não é responsável pelo animal, preocupação apontada pela maioria dos participantes da pesquisa de campo (Apêndice B, Figura~B9), o aplicativo oferece um fluxo de comprovação de tutela. Nas publicações de animais encontrados, o endereço exato e o número do imóvel de quem acolheu o animal ficam ocultos. Para liberá-los, o interessado envia uma solicitação de devolução com um relato que justifique a tutela e fotografias que possam comprová-la, como imagens antigas do animal ou documentos, como a carteira de vacinação.

A solicitação fica pendente e o autor da publicação é avisado por notificação e por mensagem privada. Ele então analisa o relato e as imagens e decide aprovar ou recusar o pedido, podendo registrar o motivo da recusa; o solicitante é notificado da decisão. Com a aprovação, o endereço exato é liberado ao solicitante para que as partes combinem a devolução. Assim, o recurso concilia a privacidade de quem acolheu o animal com a necessidade de devolvê-lo ao tutor. Até que essa ação seja formalizada pelo autor, a publicação permanece ativa no mapa e no mural.

O Capítulo 6 apresenta a arquitetura do sistema, as tecnologias empregadas e o modelo de dados que sustentam essas funcionalidades.
